\pagestyle{plain}

\chapter{CAPÍTULO 3: Adaptação Transcultural do Questionário WOCAT-SLM e Elicitação Estruturada de Saberes Agroecológicos Tradicionais via Delphi em Comunidades Quilombolas do Semiárido Nordeste II}


\begin{flushright}
Catuxe Varjão de Santana Oliveira  \footnote[1]{Programa de Pós-graduação em Ciências da Propriedade Intelectual-PPGPI/Universidade Federal de Sergipe-UFS. *E-mail:  \href{mailto:catuxe@academico.ufs.br}{\color{blue}{\underline{catuxe@academico.ufs.br}}}}
\\[0.3cm]
Luiz Diego Vidal Santos \footnote[2]{Universidade Estadual de Feira de Santana - UEFS, Brasil. E-mail: \href{mailto:ldvsantos@uefs.br}{\color{blue}{\underline{ldvsantos@uefs.br}}}}
\\[0.3cm]
Paulo Roberto Gagliardi \footnotemark[1]
\end{flushright}

% resumo em português — bookmark suprimido para manter hierarquia do capítulo
{\renewcommand{\PRIVATEbookmarkthis}[1]{}%
\begin{resumoumacoluna}
A utilização de instrumentos internacionais padronizados para avaliação de práticas de manejo sustentável da terra em contextos culturais distintos exige adaptação transcultural rigorosa, seguida de elicitação estruturada de variáveis que converta saberes tácitos em dimensões mensuráveis. O questionário WOCAT (\textit{World Overview of Conservation Approaches and Technologies}), referência global para documentação de tecnologias de gestão sustentável da terra, não dispõe de versão adaptada para comunidades tradicionais brasileiras. Este estudo integra duas fases complementares, abrangendo a adaptação transcultural do módulo de avaliação de Tecnologias SLM do questionário WOCAT para aplicação em comunidades quilombolas do Território de Identidade Semiárido Nordeste II (Bahia), conforme protocolo de \citeonline{Beaton2000} complementado pelas diretrizes da International Test Commission (ITC), e a elicitação estruturada de variáveis linguísticas para salvaguarda biocultural de Saberes e Sistemas Agrícolas Tradicionais (SSAT) mediante protocolo Delphi de três rodadas com painel heterogêneo de 15--25 especialistas. A Fase~1 envolveu seis etapas (tradução direta, síntese, retrotradução, comitê de especialistas incluindo mestres de saberes quilombolas, pré-teste com 30 agricultores em Jeremoabo (BA) e debriefing cognitivo), produzindo a versão WOCAT-SLM-QBR com IVC global $\geq$~0,90, kappa de Fleiss $\geq$~0,75 e taxa de compreensão $\geq$~85\%. O processo identificou dimensões culturalmente específicas, notadamente a espiritual-ritual e a transmissão intergeracional via oralidade, incorporadas como itens suplementares. A Fase~2, ancorada no instrumento adaptado, elicitou variáveis candidatas à incorporação no Índice de Salvaguarda Biocultural (ISB) em seis dimensões pré-estruturadas derivadas do WOCAT. A convergência foi aferida por mediana, intervalo interquartil, coeficiente de concordância de Kendall ($W$) e Coeficiente de Validade de Conteúdo (CVC). Entrevistas semiestruturadas com agricultores quilombolas triangularam o consenso do painel com a percepção êmica dos detentores dos saberes. Os resultados demonstram que o protocolo combinado adaptação--Delphi produz conjunto validado de variáveis linguísticas com CVC~$\geq$~0,80 e $W \geq$~0,70, assegurando simultaneamente equivalência transcultural e validade ecológica. O estudo contribui para a operacionalização de frameworks de salvaguarda biocultural ao fornecer cadeia metodológica integrada (da adaptação de instrumentos à elicitação consensual) que conecta a riqueza fenomenológica dos saberes tradicionais à exigência de formalização de sistemas de apoio à decisão.

\vspace{\onelineskip}
\noindent
\textbf{Palavras-chave}: Adaptação transcultural; Método Delphi; WOCAT; Gestão sustentável da terra; Elicitação de variáveis; Comunidades quilombolas; Bioeconomia; Propriedade intelectual coletiva.
\end{resumoumacoluna}}

\newpage

\section{Introdução}

A salvaguarda de Saberes e Sistemas Agrícolas Tradicionais (SSAT) em contextos de vulnerabilidade biocultural demanda uma cadeia metodológica com duas etapas estruturantes, a saber, a disponibilidade de instrumento de avaliação culturalmente equivalente, capaz de capturar as múltiplas dimensões dos sistemas produtivos tradicionais, e a conversão de conhecimentos tácitos em variáveis mensuráveis, consensuadas e auditáveis, passíveis de integração em modelos de priorização de salvaguarda e sistemas de apoio à decisão \cite{Nonaka1995,CohenLevinthal1990}. O presente estudo integra essas duas etapas em protocolo sequencial articulado.

O questionário WOCAT (\textit{World Overview of Conservation Approaches and Technologies}), desenvolvido pelo Centre for Development and Environment (CDE) da Universidade de Berna e adotado como ferramenta oficial pela FAO, UNCCD e rede global de parceiros \cite{Liniger2019,Schwilch2011}, representa o framework mais consolidado para documentação de tecnologias SLM. Com aplicações em mais de 120 países e banco de dados contendo mais de 2.000 tecnologias catalogadas, o WOCAT oferece arquitetura padronizada de sete seções que abrangem desde a classificação técnica da tecnologia até a avaliação de impactos socioeconômicos, ecológicos e adaptativos, passando por análise de custos, ambiente biofísico e governança dos usuários da terra.

Contudo, a aplicação direta desse instrumento a comunidades tradicionais brasileiras, particularmente comunidades quilombolas do semiárido nordestino, esbarra em pelo menos três barreiras de equivalência. Primeiro, a \textbf{barreira semântico-idiomática}: o questionário original em inglês contém terminologia técnica agronômica sem tradução direta para o português vernacular de agricultores quilombolas, cujo vocabulário opera sob categorias etnotaxonômicas próprias \cite{Rist2006,Quave2014}. Segundo, a \textbf{barreira experiencial}: categorias de resposta pressupõem contextos fundiários formalizados, enquanto comunidades quilombolas operam sob regimes coletivos de uso com reconhecimento jurídico precário \cite{Almeida2011}. Terceiro, a \textbf{barreira conceitual}: o WOCAT organiza-se sob racionalidade agronômica ocidental que não contempla dimensões espirituais, rituais e simbólicas constitutivas da lógica de manejo quilombola \cite{Santos2007,Toledo2008}.

Essas barreiras não são meramente acadêmicas. A literatura sobre equivalência transcultural \cite{Herdman1999,Guillemin1993} demonstra que a aplicação de instrumentos sem adaptação formal produz viés sistemático de mensuração. No contexto desta tese, esse risco assume magnitude crítica: o WOCAT-SLM servirá como template de referência para a elicitação Delphi e, indiretamente, para a construção do Índice de Salvaguarda Biocultural (ISB) via integração TRI-fuzzy (OE3). Um instrumento mal adaptado propagaria erros sistêmicos por toda a cadeia.

Paralelamente, os SSAT permanecem à margem dos mecanismos formais de salvaguarda e proteção, a despeito de sua relevância para a segurança alimentar e a resiliência territorial \cite{Altieri1995,Toledo2008}. Essa invisibilidade decorre de lacuna metodológica: a ausência de protocolos validados para converter saberes tácitos, oralmente transmitidos e corporalmente incorporados \cite{Polanyi1966}, em variáveis mensuráveis e auditáveis que possam instruir políticas de salvaguarda e sistemas de apoio à decisão.

Conforme demonstrado pela revisão de escopo PRISMA-ScR conduzida no Capítulo 1, embora o Aprendizado de Máquina alcance acurácias combinadas superiores a 90\% no monitoramento de paisagens agrícolas, a conformidade FAIR permanece criticamente baixa e a ausência de protocolos de validação espacial compromete a prontidão operacional. Antes de se aplicar qualquer tecnologia computacional sofisticada, é imprescindível dispor de variáveis de entrada estruturadas, validadas e contextualmente pertinentes.

Nesse cenário, o método Delphi \cite{Linstone1975,Rowe1999} emerge como técnica de elicitação particularmente adequada por operar sob anonimato e feedback controlado (minimizando efeitos de dominância social), por permitir convergência mensurável mediante coeficientes estatísticos \cite{Diamond2014}, e por assegurar validade ecológica quando operacionalizado com painéis heterogêneos que incluem detentores de saberes tradicionais \cite{Santos2007}. Aplicado sobre o instrumento previamente adaptado transculturalmente, o Delphi produz variáveis com dupla legitimidade: técnica (consenso estatístico interavaliadores) e cultural (ancoragem no WOCAT-SLM-QBR validado junto às comunidades).

As questões norteadoras (Q2) integram ambas as fases: \textit{Em que medida o processo combinado de adaptação transcultural e elicitação Delphi produz instrumento culturalmente equivalente e conjunto validado de variáveis mensuráveis, aumentando consenso e comparabilidade entre avaliadores?} A hipótese (H2) postula que: (a)~a versão adaptada (WOCAT-SLM-QBR) atingirá IVC~$\geq$~0,80, kappa de Fleiss~$\geq$~0,70 e taxa de compreensão~$\geq$~80\%, gerando itens suplementares culturalmente específicos; e (b)~a aplicação do método Delphi sobre as dimensões do instrumento adaptado produzirá índice de consenso estatístico ($W \geq 0,70$; $CVC \geq 0,80$) superior ao obtido por métodos não estruturados.

O estudo insere-se no Objetivo Específico 2 (OE2) de um programa de pesquisa dedicado ao desenvolvimento de sistema computacional de salvaguarda biocultural para o Território de Identidade Semiárido Nordeste II (Bahia), articulando revisão de escopo (OE1, Capítulo~1), meta-análise de vulnerabilidade (Capítulo~2), adaptação transcultural e elicitação estruturada (OE2, este Capítulo), modelagem psicométrica e fuzzy (OE3, Capítulo~4), auditoria computacional (OE4, Capítulo~5) e design participativo (OE5, Capítulo~6). Nessa arquitetura, o presente capítulo opera como ponte metodológica entre o diagnóstico tecnológico do Capítulo~1 e a construção do motor de inferência fuzzy, fornecendo simultaneamente instrumento culturalmente calibrado e variáveis de entrada validadas.


\section{Referencial Teórico}

\subsection{Teoria da Equivalência Transcultural de Instrumentos}

A utilização de instrumentos desenvolvidos em um contexto cultural para aplicação em outro exige mais do que tradução linguística. \citeonline{Herdman1999} propuseram modelo hierárquico de equivalência transcultural composto por seis níveis: equivalência conceitual, de itens, semântica, operacional, de mensuração e funcional.

A distinção entre abordagens \textit{etic} (universalista) e \textit{emic} (culturalmente específica) \cite{Pike1967} é particularmente relevante. O WOCAT adota perspectiva predominantemente \textit{etic}, pressupondo categorias universalmente aplicáveis. Embora essa abordagem viabilize comparabilidade internacional, pode obscurecer categorias \textit{emic} significativas. Agricultores quilombolas classificam terras por atributos espirituais ou memória social, categorias invisíveis ao WOCAT original \cite{Toledo2008}. A adaptação transcultural visa preservar a dimensão \textit{etic} que confere comparabilidade enquanto incorpora dimensões \textit{emic} que conferem validade ecológica.

\citeonline{Guillemin1993} formalizaram o protocolo de adaptação refinado por \citeonline{Beaton2000}. Os princípios fundamentais compreendem tradução por mais de um tradutor independente, retrotradução por tradutores não familiarizados com o original, revisão por comitê multidisciplinar e multicultural, e pré-teste com amostra da população-alvo.

\subsection{Conhecimento Tácito, Espiral SECI e Conversão de Saberes}

A distinção entre conhecimento tácito e explícito de \citeonline{Polanyi1966} constitui ponto de partida epistemológico desta investigação. Nos sistemas agroecológicos tradicionais, a dimensão tácita manifesta-se com particular intensidade: manejo fenológico, leitura de sinais climáticos, seleção de variedades adaptadas e práticas de conservação de solos são transmitidos oralmente e pela prática cotidiana, resistindo à codificação em categorias discretas \cite{Berkes2017,Toledo2008}.

A espiral SECI de \citeonline{Nonaka1995} (Socialização, Externalização, Combinação e Internalização) fornece arcabouço conceitual para compreender como o conhecimento tácito pode ser convertido progressivamente. No contexto deste estudo, a adaptação transcultural opera na transição Socialização--Externalização (traduzindo práticas locais em linguagem do instrumento), enquanto o Delphi opera na fase de \textit{externalização} avançada, onde julgamentos qualitativos são traduzidos em variáveis linguísticas calibradas e escalas padronizadas. Essa externalização estruturada distingue-se criticamente da espontânea por submeter o resultado a critérios de consenso verificáveis.

\subsection{Método Delphi: Fundamentos e Evolução}

O método Delphi, desenvolvido pela RAND Corporation \cite{Linstone1975}, evoluiu para ferramenta consolidada de construção de consenso em domínios onde dados empíricos são escassos. \citeonline{Rowe1999} identificaram quatro características definidoras: anonimato, iteração com feedback controlado, agregação estatística e heterogeneidade controlada do painel.

No contexto de conhecimentos tradicionais, o anonimato é particularmente relevante para neutralizar assimetrias de poder entre especialistas acadêmicos e detentores de saberes locais \cite{Hasson2000}. A mensuração do consenso utiliza múltiplos indicadores: coeficiente de concordância de Kendall ($W \geq 0,70$ como consenso forte) \cite{Schmidt2014}, intervalo interquartil ($IQR \leq 1$ em escala de 5 pontos como limiar operacional) \cite{VonDerGracht2012} e Coeficiente de Validade de Conteúdo ($CVC \geq 0,80$) \cite{Hernandez-Nieto2002}.

A articulação entre adaptação transcultural e Delphi constitui contribuição metodológica original: o instrumento adaptado provê template estruturado de dimensões validadas culturalmente, sobre o qual o Delphi opera para elicitar, refinar e consensuar variáveis específicas. Essa sequência evita tanto a imposição de categorias \textit{ad hoc} quanto a dispersão inerente a processos de elicitação sem ancoragem instrumental.

\subsection{O Framework WOCAT e Operacionalização das Dimensões}

O questionário WOCAT para Tecnologias SLM apresenta arquitetura modular em sete seções \cite{Liniger2019}, das quais cinco fornecem conteúdo operacionalizável para as dimensões de vulnerabilidade biocultural ($V_1$--$V_6$) adotadas nesta tese. A seção~2 (Descrição da Tecnologia SLM), que abrange definição, especificações técnicas e classificação de medidas, alimenta $V_1$ (transmissão intergeracional) e $V_4$ (completude de registros). A seção~3 (Classificação da Tecnologia), relativa a uso da terra, degradação e função protetora, fornece indicadores para $V_2$ (complexidade biocultural) e $V_3$ (singularidade territorial). A seção~5 (Ambiente Natural e Humano), que cobre clima, topografia, solos e características dos usuários, contribui para $V_2$, $V_3$ e $V_5$ (regime fundiário e direitos de uso). A seção~6 (Impactos e Conclusões), com avaliação de impactos socioculturais, ecológicos e institucionais, sustenta $V_1$, $V_2$, $V_4$ e $V_6$ (organização social). A seção~7 (Referências e Links) complementa $V_4$ ao indicar a existência de publicações e registros formais. As seções~1 (Informações Gerais, de natureza cadastral) e~4 (Insumos e Custos) não geram variáveis de vulnerabilidade independentes e serão tratadas como contexto descritivo na ficha de campo.

Para a adaptação transcultural (Fase~1), o foco recai sobre as seções 2, 3, 5 e 6, que contêm os construtos avaliativos culturalmente sensíveis. Para a elicitação Delphi (Fase~2), essas mesmas seções, acrescidas da seção~7, informam as dimensões pré-estruturadas do painel, conforme mapeamento na Tabela~\ref{tab:wocat_delphi}. Essa correspondência explicita como os módulos do instrumento WOCAT se convertem em variáveis operacionalizáveis dentro do framework hexadimensional da tese, evitando fragmentação taxonômica e inflação de parâmetros.

\begin{table}[htbp]
\centering
\caption{Mapeamento entre seções do questionário WOCAT, dimensões de vulnerabilidade biocultural ($V_1$--$V_6$) e variáveis-chave deriváveis para o painel Delphi.}
\label{tab:wocat_delphi}
\small
\begin{tabular}{p{1cm}p{2.8cm}p{4cm}p{5cm}}
\toprule
\textbf{Dim.} & \textbf{Denominação} & \textbf{Seções WOCAT de origem} & \textbf{Variáveis-chave deriváveis} \\
\midrule
$V_1$ & Erosão Intergeracional & §2 Descrição; §6.1 Impactos socioculturais & Proporção de jovens praticantes, taxa de transmissão ativa, significado ritual vinculado à prática, oportunidades culturais de aprendizado \\
$V_2$ & Complexidade Biocultural & §3 Classificação; §5 Ambiente natural; §6.1 Impactos ecológicos & Agrobiodiversidade (Shannon $H'$, Simpson), riqueza de variedades locais, interações solo-planta codificadas, resiliência edáfica \\
$V_3$ & Singularidade Territorial & §3 Classificação; §5 Ambiente natural & Beta-diversidade de práticas entre comunidades, endemismo de variedades, proporção de práticas exclusivas \\
$V_4$ & Status de Documentação & §2 Descrição; §6.5 Adoção; §7 Referências & Número de saberes com registro formal, completude de fichas técnicas, taxa de digitalização de acervos \\
$V_5$ & Vulnerabilidade Jurídica & §5.6 Características dos usuários; §5.8 Propriedade & Mecanismos formais de proteção (IG, marca coletiva, registro IPHAN), existência de protocolo comunitário, regime fundiário, direitos de uso \\
$V_6$ & Organização Social & §5.9 Infraestrutura; §6.1 Instituições comunitárias; §6.5 Adoção & Mestres de saberes ativos, frequência de eventos de transmissão, redes de cooperação, capital social, equidade de gênero \\
\bottomrule
\end{tabular}
\end{table}

Conteúdos avaliativos do WOCAT que não geram dimensão autônoma (§4, Insumos e custos, e indicadores adaptativos de §3.8 e §6.3) foram absorvidos como variáveis transversais em $V_2$ (resiliência edáfica), $V_3$ (exclusividade de práticas adaptativas) e $V_6$ (capacidade organizacional de resposta a perturbações), preservando a parcimônia do modelo hexadimensional validado pela meta-análise (Capítulo~2).

\subsection{Comunidades Quilombolas e Especificidades Ontológicas dos SSAT}

Os SSAT quilombolas operam sob lógica que difere estruturalmente da racionalidade agronômica convencional. \citeonline{Toledo2008} demonstram que sistemas tradicionais são governados pelo complexo Conhecimento-Prática-Crença (K-P-B), onde crenças funcionam como regulador ético-cosmológico do manejo. \citeonline{Berkes2017} complementa que esses sistemas constituem exemplos de manejo adaptativo onde experimentação empírica é modulada por instituições sociais.

No contexto quilombola do semiárido, pelo menos três dimensões não são capturadas pelo WOCAT original. A dimensão espiritual-ritual, que abrange bênçãos sobre sementes, rituais de plantio conforme ciclos lunares e proibições em datas sagradas, integra o sistema produtivo mas permanece invisível ao WOCAT. A dimensão de transmissão intergeracional via oralidade também não é contemplada, uma vez que o WOCAT documenta a tecnologia como produto acabado sem capturar o processo de transmissão oral que constitui mecanismo de inovação dos SSAT \cite{Polanyi1966}. De modo análogo, a dimensão coletiva-comunitária fica subrepresentada, dado que muitas práticas quilombolas são coletivas (mutirões, trocas de sementes, manejo comunitário) e operam sob lógica de bens comuns \cite{Ostrom1990}.

O Território de Identidade Semiárido Nordeste II (Bahia) compreende 18 municípios. Jeremoabo concentra 11 comunidades quilombolas certificadas pela Fundação Cultural Palmares, que mantêm sistemas agroflorestais e práticas de manejo adaptados ao semiárido \cite{Altieri1995}. Essas comunidades enfrentam duplo desafio: erosão acelerada dos saberes sob pressão de monoculturas e inexistência de instrumentos formais para traduzir a sofisticação de seus sistemas em linguagem acessível a mercados e marcos regulatórios.

\subsection{Gestão do Conhecimento, ISO~30401 e Rastreabilidade}

A conversão de saberes tácitos em variáveis mensuráveis exige governança do conhecimento que assegure rastreabilidade, retenção e transferibilidade. A ISO~30401:2018 (Sistemas de Gestão do Conhecimento) reconhece que o valor do conhecimento depende de cultura, processos e aprendizagem, não apenas de codificação estática \cite{ISO30401}. Em comunidades quilombolas onde o conhecimento tácito está concentrado em poucos mestres de saberes, a documentação via instrumento adaptado e protocolo Delphi constitui exercício urgente de preservação.

A integração entre adaptação transcultural, Delphi e ISO~30401 materializa-se em três frentes complementares. Cada variável validada é acompanhada de definição operacional com origem rastreável, enquanto o dossiê de adaptação documenta todo o processo de externalização. A devolutiva às comunidades fecha o ciclo de internalização (SECI).

Essa perspectiva dialoga com a \textit{capacidade absortiva reversa}: são as instituições formais que precisam adaptar suas ferramentas para absorver o conhecimento das comunidades \cite{CohenLevinthal1990}. O instrumento adaptado materializa essa inversão, na qual não são os quilombolas que se ajustam ao WOCAT, mas o WOCAT que se adapta à realidade quilombola \cite{Santos2007}.


\section{Materiais e Métodos}

\subsection{Delineamento Geral}

Esta investigação configura-se como estudo metodológico de métodos mistos \cite{Creswell2018}, organizado em duas fases sequenciais integradas: Fase~1, Adaptação Transcultural (pesquisa metodológica, protocolo de \citeonline{Beaton2000}), e Fase~2, Elicitação Delphi (estudo exploratório-descritivo com triangulação qualitativa). A integração segue delineamento sequencial: o instrumento adaptado na Fase~1 serve como template de referência para as dimensões avaliativas da Fase~2.

% ============================================================
% FASE 1: ADAPTAÇÃO TRANSCULTURAL
% ============================================================
\subsection{Fase 1: Adaptação Transcultural do WOCAT-SLM}

\subsubsection{Etapa 1: Tradução Direta (Inglês $\rightarrow$ Português)}

Dois tradutores independentes realizarão tradução do questionário WOCAT-SLM (seções 2, 3, 5 e 6) do inglês para o português brasileiro. O primeiro (T1), profissional com formação em ciências agrárias, bilíngue e ciente dos objetivos, priorizará equivalência técnica. O segundo (T2), profissional sem formação técnica na área, bilíngue e não informado dos objetivos, priorizará linguagem coloquial.

A divergência intencional entre perfis maximiza a detecção de ambiguidades \cite{Beaton2000}. Cada tradutor produzirá versão independente (T1 e T2) com relatório de decisões.

\subsubsection{Etapa 2: Síntese das Traduções (T-12)}

Os tradutores e um mediador produzirão versão sintetizada (T-12). Discrepâncias serão resolvidas mediante negociação documentada, com itens não resolvidos encaminhados ao comitê de especialistas.

\subsubsection{Etapa 3: Retrotradução (Português $\rightarrow$ Inglês)}

Dois retrotradutores independentes, nativos de língua inglesa ou com proficiência C2, sem conhecimento do original, traduzirão a T-12 de volta para o inglês (BT1 e BT2). As retrotraduções serão comparadas com o instrumento original item a item.

\subsubsection{Etapa 4: Comitê de Especialistas}

O comitê multidisciplinar será composto por 8 a 12 membros, incluindo mestres de saberes quilombolas (2 a 3 agricultores com experiência mínima de 20 anos), pesquisadores em agroecologia e etnoecologia (2 a 3 doutores com experiência participativa), especialistas em psicometria e adaptação transcultural (1 a 2) e gestores de PI e extensionistas (2 a 3).

A inclusão de mestres de saberes como membros plenos fundamenta-se no princípio de soberania epistêmica \cite{Santos2007} e nas diretrizes ITC \cite{ITC2017}.

O comitê avaliará cada item em quatro dimensões de equivalência (escala de 4 pontos): semântica, idiomática, experiencial e conceitual. O Índice de Validade de Conteúdo será calculado:

\begin{equation}
IVC_{item} = \frac{\text{nº de avaliadores que atribuíram 3 ou 4}}{\text{nº total de avaliadores}}
\end{equation}

Critérios: $IVC \geq 0,80$ (aceito); $0,60 \leq IVC < 0,80$ (revisão); $IVC < 0,60$ (reformulação/exclusão). Concordância interavaliadores aferida por kappa de Fleiss \cite{Fleiss1971}. O comitê também identificará lacunas culturais e proporá itens suplementares.

\subsubsection{Etapa 5: Pré-Teste}

Versão pré-final aplicada a 30--40 agricultores quilombolas de Jeremoabo (BA), selecionados por amostragem intencional com variabilidade em idade, gênero, escolaridade e sistema produtivo. Aplicação em formato de entrevista assistida. Após cada seção, \textbf{debriefing cognitivo} \cite{Willis2005} com perguntas padronizadas de compreensão, alternativas linguísticas e pertinência experiencial.

Indicadores: taxa de compreensão por item ($\geq$ 85\%), taxa de não-resposta ($\leq$ 15\%), tempo de aplicação e distribuição das respostas (efeito teto/piso).

\subsubsection{Etapa 6: Consolidação}

Versão final WOCAT-SLM-QBR consolidada com dossiê completo de adaptação: versões T1, T2, T-12, BT1, BT2, atas do comitê, dados do pré-teste, manual de aplicação e versão final. Encaminhamento ao WOCAT Secretariat.

% ============================================================
% FASE 2: ELICITAÇÃO DELPHI
% ============================================================
\subsection{Fase 2: Elicitação Estruturada via Protocolo Delphi}

\subsubsection{Composição do Painel}

O painel será composto por 15 a 25 participantes selecionados por amostragem intencional \cite{Patton2015}, buscando heterogeneidade controlada em quatro categorias: mestres de saberes quilombolas (4 a 6 agricultores reconhecidos por suas comunidades, com experiência mínima de 20 anos), pesquisadores acadêmicos (4 a 6 doutores em agroecologia, etnoecologia, PI ou gestão da inovação), técnicos extensionistas (3 a 5, com experiência mínima de 5 anos em assessoria a comunidades tradicionais) e gestores de PI e bioeconomia (3 a 5, vinculados a NITs, SEBRAE, INPI ou secretarias territoriais).

Critérios: experiência mínima de 5 anos, reconhecimento pela comunidade epistêmica/territorial e disponibilidade para 3 rodadas em 4 meses. A inclusão de mestres de saberes como especialistas de pleno direito é fundamentada em \cite{Santos2007,Beaton2000}.

\subsubsection{Estrutura das Rodadas}

\textbf{Rodada 1, Divergência e Exploração.} Questionário aberto solicitando enumeração de variáveis relevantes para avaliação de vulnerabilidade e salvaguarda de SSAT, organizadas nas seis dimensões de vulnerabilidade biocultural ($V_1$--$V_6$) conforme mapeamento WOCAT (Tabela~\ref{tab:wocat_delphi}). Mestres de saberes com preferência por participação oral terão respostas transcritas por facilitadores. Consolidação mediante análise de conteúdo \cite{Braun2006}.

\textbf{Rodada 2, Convergência.} Questionário estruturado com lista consolidada. Avaliação em escala Likert de 5 pontos para três critérios: relevância, clareza e operacionalidade. Cálculo de mediana ($Md$), intervalo interquartil ($IQR$), coeficiente de variação ($CV$) e frequência de respostas extremas.

\textbf{Rodada 3, Consenso.} Reenvio com feedback agregado (medianas, distribuição, posicionamento individual anonimizado). Consenso operacional: $IQR \leq 1,0$ e $Md \geq 4,0$. Variáveis sem consenso submetidas a análise qualitativa complementar.

\subsubsection{Análise Estatística do Consenso}

A análise estatística do consenso envolverá o cálculo de mediana e IQR por variável e por rodada, avaliando-se a redução progressiva do IQR. O coeficiente de concordância de Kendall ($W$) será empregado como indicador de consenso, adotando $W \geq 0{,}70$ como consenso forte e $0{,}50 \leq W < 0{,}70$ como moderado \cite{Schmidt2014}. O Coeficiente de Validade de Conteúdo (CVC) seguirá critério $CVC \geq 0{,}80$ \cite{Hernandez-Nieto2002}. A taxa de estabilidade entre as rodadas 2 e 3 será aferida pelo percentual de painelistas com ajuste $\pm 1$ ponto, e o teste de Friedman ($\alpha = 0{,}05$) com \textit{post-hoc} de Dunn será aplicado quando aplicável.

Para teste direto da H2(b), comparação entre índices de consenso Delphi e levantamento qualitativo não estruturado (grupo focal com 8--10 especialistas do mesmo universo). Diferença testada por permutação ou bootstrap.

\subsection{Componente de Triangulação: Entrevistas Semiestruturadas}

Entrevistas com 15--20 agricultores quilombolas de Jeremoabo, selecionados por saturação teórica \cite{Glaser1967}. Roteiro abordando: percepção sobre variáveis do Delphi, dimensões não contempladas, adequação da linguagem e hierarquização espontânea.

Gravação em áudio, transcrição integral e análise temática \cite{Braun2006} (familiarização, codificação aberta, busca por temas, revisão, redação). Codificação por dois pesquisadores independentes (kappa de Cohen $\geq 0,60$). Triangulação via matriz de correspondência variáveis validadas $\times$ categorias temáticas emergentes.

\subsection{Aspectos Éticos}

O projeto será submetido ao Comitê de Ética em Pesquisa da UFS (Resolução CNS nº~466/2012 e nº~510/2016). Consentimento livre, prévio e informado de todos os participantes. Mestres de saberes do comitê reconhecidos como coautores do instrumento. Dados sensíveis tratados conforme Protocolo de Nagoia e Lei nº~13.123/2015, com devolutiva às comunidades.


\section{Resultados Esperados}

Espera-se como produtos diretos desta investigação a versão WOCAT-SLM-QBR, instrumento adaptado transculturalmente para contexto quilombola brasileiro contendo itens traduzidos e adaptados acrescidos de itens suplementares culturalmente específicos (IVC global $\geq$~0,90 e kappa~$\geq$~0,75), acompanhada de dossiê de adaptação com documentação auditável de todas as decisões, constituindo referência para adaptações em outros contextos de comunidades tradicionais brasileiras (indígenas, ribeirinhas, caiçaras, fundo de pasto). O estudo deverá produzir ainda um mapa de lacunas culturais, identificando as dimensões não contempladas pelo WOCAT original (espera-se ao menos as dimensões espiritual-ritual, de transmissão intergeracional por oralidade e de coletividade como bens comuns), cujas variáveis serão integradas como itens suplementares nas dimensões V1–V6 pertinentes em vez de constituírem dimensões autônomas. Prevê-se também um conjunto validado de variáveis linguísticas, estimado em 20 a 30 variáveis, acompanhadas de definições operacionais consensuadas e organizadas nas seis dimensões de vulnerabilidade biocultural ($V_1$--$V_6$) estabelecidas pela meta-análise (Capítulo~2). Prevê-se igualmente a produção de evidência empírica de superioridade do método estruturado (H2b), quantificada pela diferença entre coeficientes de concordância Delphi \textit{versus} grupo focal não estruturado, e de um manual de operacionalização do protocolo integrado adaptação-Delphi para contextos de comunidades tradicionais, documentando adaptações metodológicas relativas a facilitação oral, tradução linguística e gestão de assimetrias de poder. A matriz de correspondência entre consenso técnico (Delphi) e percepção comunitária (entrevistas) fornecerá base de validade ecológica, e as variáveis validadas constituirão insumo direto para o OE3 (Capítulo~4), operando como antecedentes do motor de inferência fuzzy do Índice de Salvaguarda Biocultural (ISB).


\section{Discussão Preliminar}

A integração entre adaptação transcultural e elicitação Delphi em protocolo sequencial articulado responde a uma lacuna identificada na literatura de salvaguarda de conhecimentos tradicionais \cite{Berkes2017,Toledo2008}: a ausência de cadeia metodológica que conecte rigor na construção de instrumentos validados com consenso auditável sobre variáveis mensuráveis, mantendo legitimidade cultural junto aos detentores dos saberes. Enquanto a adaptação transcultural assegura que o instrumento mede o que pretende medir no contexto quilombola, o Delphi fornece estrutura estatística para consenso interavaliadores, e as entrevistas garantem que esse consenso não se descole das ontologias locais de conhecimento \cite{Santos2007}.

As variáveis elicitadas operacionalizam as dimensões de vulnerabilidade biocultural ($V_1$--$V_6$), criando infraestrutura para que comunidades quilombolas disponham de diagnóstico verificável sobre o estado de seus saberes e sistemas agrícolas. Essa documentação estruturada constitui pré-requisito para negociações de repartição de benefícios, certificação de produtos e proteção jurídica via indicações geográficas ou marcas coletivas \cite{Belletti2015}.

A inclusão de mestres de saberes quilombolas como membros plenos tanto do comitê de adaptação quanto do painel Delphi representa inovação metodológica alinhada ao paradigma da soberania epistêmica \cite{Santos2007}. Os detentores de saberes tradicionais são coautores do instrumento e do consenso, exercendo agência sobre como sua realidade será representada e mensurada.

A expectativa de que o processo gere itens suplementares não contemplados pelo WOCAT original corrobora a limitação de abordagens puramente \textit{etic}. As dimensões espiritual-ritual, de oralidade e de coletividade constituem evidência de que frameworks universalistas carregam pressupostos culturais que operam como ``pontos cegos'' quando transplantados para ontologias distintas \cite{Herdman1999}.

Na arquitetura mais ampla desta tese, o presente capítulo opera como fundação metodológica: o instrumento culturalmente calibrado e as variáveis consensuadas alimentarão diretamente as funções de pertinência e regras SE-ENTÃO do sistema fuzzy Mamdani (OE3/Capítulo~4). Cada variável do ISB terá origem documentada em adaptação transcultural e consenso especializado, conferindo rastreabilidade ao modelo \cite{Costanza1997} e atendendo à cadeia de evidências: WOCAT original $\rightarrow$ WOCAT-SLM-QBR $\rightarrow$ consenso Delphi $\rightarrow$ meta-análise $\rightarrow$ TRI $\rightarrow$ modelo fuzzy.


%% Textos comentados para desenvolvimento futuro
%% Seção de Resultados será preenchida após coleta de dados de campo
%% Seção de Conclusões será redigida após análise dos resultados
